\begin{afterword}
\summaryhead

The post opens with two epigraphs from Japanese martial texts: a ``strong'' effort gives only mediocre results, and success needs an effort ``as if our life is at stake.'' It introduces isshokenmei, a Japanese word for a desperate all-out effort, and says that self-described rationalists need such an effort to stop making mistakes. Its example is Robert Aumann, who proved a theorem about Bayesian agreement and is an Orthodox Jew.

Yet Japan, whose culture honours desperate effort and says that ``The nail that sticks up gets hammered down,'' does not lead the United States in science. The West's taste for risk and for leaving the herd seems to count as much. So the author names a higher virtue: extraordinary effort, doing something out of the ordinary, as one would need to lift a truck. It is also more dangerous, and rules should be broken only for an overwhelmingly good reason.

The author's examples are people who called Friendly AI futile and people who advised earning degrees first. The post asks for a sense of the cost of ordinariness as well as of the risk of the extraordinary, and ends with ``It's all part of the plan.''

\responsehead

The post's main idea is a distinction between two kinds of effort. A desperate effort does more of the usual thing; an extraordinary effort does something different. The truck makes the difference plain. The post also weighs the danger of its own advice: its embezzler goes to extraordinary lengths too, and its maxim is to break rules only for ``an overwhelmingly good reason.'' It ends by asking for a sense of both the risk of the extraordinary and the cost of the ordinary. That is a balanced rule, and the post states it clearly.

The case that the distinction matters rests on a comparison between Japan and the United States: a Japanese word, two proverbs, the seriousness of Japanese exams, and the fact that Japan does not lead in science. The post hedges the inference (``this would seem to count''). Still, one pair of countries cannot separate these traits from the many other ways the two countries differ. The later sentence that extraordinary effort is not included in isshokenmei, ``or Japan would be a very different place,'' rests on the same comparison.

The illustrations of the extraordinary are of two kinds. Two are warnings: the embezzler and the friend who wanted to change a game's rules for no reason. The other two come from the author's own work: people who called Friendly AI futile, and people who advised degrees first. Whether the extraordinary course was right in these cases depends on a premise that the post states and does not argue: that ``the life or death of the human species'' depends on a few people doing something a little extraordinary. The post says so itself; it sets ``the details of that debate'' aside. The reader is shown the attitude and given a test for it, but not a case in which the test was applied and the extraordinary course turned out right.

Two smaller points are in the notes. The Aumann example repeats earlier posts; the notes on ``Crisis of Faith'' give Aumann's own account of his religion. And ``Crisis of Faith,'' placed earlier in the book but written three days later, uses isshokenmei for ``the desperate, extraordinary, convulsive effort,'' joining the two virtues this post separates.

In short: a clear distinction between desperate and extraordinary effort, with its own warnings attached, supported by one comparison between Japan and the United States and illustrated mainly by the author's own plans, whose premise the post sets aside.
\end{afterword}
